\documentclass[a4paper]{article}     
\usepackage{amsfonts}
\usepackage{amsmath}
\usepackage{amssymb}
\usepackage{graphicx}
\usepackage{cancel}

\newcommand{\asa}{\Leftrightarrow} %als slechts als
\newcommand{\adR}{\Rightarrow} %als dan Rechts
\newcommand{\adL}{\Leftarrow}
\newcommand{\Lh}{\left(} %Links haakje
\newcommand{\Rh}{\right)}
\newcommand{\Lv}{\left[} %Links vierkanthaakje
\newcommand{\Rv}{\right]}
\newcommand{\La}{\left\{} %Links acolade
\newcommand{\Ra}{\right\}}
\newcommand{\Ras}{\right.} %Rechts acolade stelsel (omdat om stelsels te maken heb je een punt nodig
\newcommand{\pnR}{\rightarrow} %pijl naar Rechts
\newcommand{\limiet}{\lim\limits_}
\newcommand{\schrap}{\cancel}
\newcommand{\schrapb}{\bcancel} %backslash
\newcommand{\schrapk}{\xcancel} %kruis
\newcommand{\vervang}{\cancelto} 

\begin{document}

\noindent \large Auteur: Yoren Collier \\
\noindent \large Studierichting: Informatica\\
\noindent \large Oefeningenreeks 2D nr. 5.6\\

\medskip

\normalsize

$f(x) = \dfrac{x}{2x+1}$\\

$\limiet{x \pnR 5} f(x) = \dfrac{5}{10+1} = \dfrac{5}{11}$\\

$\limiet{x \pnR \tfrac{-1}{2}} f(x) = \dfrac{\dfrac{-1}{2}}{2\Lh \dfrac{-1}{2} \Rh + 1} = \dfrac{\dfrac{-1}{2}}{0}$\\

$\adR \limiet{x \pnR \tfrac{-1}{2}-} f(x) = +\infty \text{  en } \limiet{x \pnR \tfrac{-1}{2}+} f(x) = -\infty$

\begin{thebibliography}{999}
%\bibitem{a} Referentie
\end{thebibliography}
\end{document} 